% EXTRACTED FOR PAPER TWO (the hardware case study) -- PASS 48.1
% source     : paper/paper.tex
% blob       : 13590dbfe8119022ae95ff7a374910e30146ab8c
% lines      : 1075-1083 (1-indexed, inclusive)
% section    : Section 5.4
% prose words: 83
%
% The 'entangling overhead is real but not the binding constraint' paragraph and its forward reference to the hardware section. The claim itself survives in paper one on the transpilation counts; the forward reference does not.
%
% This is a VERBATIM COPY. The original is unmodified and still frozen.
% Regenerate with: experiments/pass48_extract_paper2.py

This speaks to the resource-cost concern behind the choice in
\cite{qmladvantage2026} to keep its measure-first protocol single-copy.
The entangling overhead that motivates such caution is real, and it is not
the binding constraint. The binding constraint is that at the tested budgets
the local-randomized-measurement shadow U-statistic we study becomes
statistically unusable for moment estimation at comparable system sizes,
a statement about this protocol, not about every single-copy
strategy. Whether the collective route survives on a real device is the
question we take up next.
